\chapter{Discussion, Screenshots, and User Interfaces}
\label{ch:discussion_ui}

\section{Prototype User Interfaces and Workflow Walkthrough}
The Track frontend delivers an intuitive, responsive user experience optimized for desktop and mobile displays. The interface comprises four primary views:

\subsection{Interactive Application Kanban Board}
The central dashboard displays candidate applications grouped into columns corresponding to the lifecycle states: \textbf{Saved}, \textbf{Applied}, \textbf{Screening}, \textbf{Interviewing}, \textbf{Offer}, and \textbf{Rejected}.

Each application card renders:
\begin{itemize}
    \item Company name, job title, and salary range (if captured).
    \item Days spent in current stage (computed dynamically to track application velocity).
    \item A visual \textbf{Stage Lock} indicator badge (depicting whether automated email workers are permitted to alter this application's status).
    \item Direct action menus to trigger stage promotion, stage lock toggling, or view the complete audit timeline.
\end{itemize}

\subsection{Job Discovery and Career Portal Management}
The discovery portal provides two views:
\begin{itemize}
    \item \textbf{Aggregated Job Stream}: A searchable, paginated catalog of all ingested jobs with real-time keyword search, remote/hybrid filters, and a single-click ``Track Application'' button. Duplicate postings are automatically consolidated with multi-source attribution badges.
    \item \textbf{Career Pages Admin Console}: An administrative interface allowing users to register target company career boards by specifying provider type (\textit{Greenhouse, Lever, Ashby, Workable}) and company slug, displaying real-time sync telemetry (\textit{Total Jobs Found}, \textit{Last Synced At}, \textit{Sync Status}).
\end{itemize}

\subsection{Recruiter Mailbox Review and Intent Classification}
The Mailbox Review screen surfaces emails synchronized through the Google OAuth2 Pub/Sub integration. When an email indicates an interview request or rejection, a review modal highlights the extracted company name, role, and suggested transition alongside current application data, allowing the user to confirm or decline the proposed transition with zero cognitive friction.

\subsection{AI Career Copilot Drawer}
Accessible from anywhere in the application, the AI Copilot slide-over drawer features:
\begin{itemize}
    \item \textbf{ATS Keyword Heatmap}: Visualizing high-frequency missing keywords between candidate resume text and target job specifications.
    \item \textbf{Google XYZ Bullet Rewriter}: An interactive form that prompts for raw project outcomes and generates structured achievement bullets formatted according to Google's standard recruitment rubric (\textit{Accomplished [X] as measured by [Y] by doing [Z]}).
\end{itemize}

\section{Engineering Challenges Overcome}

\subsection{Resolving the Destructive Overwrite Problem}
A significant challenge in building autonomous email tracking systems is the risk of conflicting status transitions. For instance, an applicant may interview for a Senior Engineer position while simultaneously receiving an automated rejection email for an older entry-level requisition at the same enterprise. Without strict safeguards, an automated email listener would parse the incoming rejection and immediately downgrade the applicant's status to \texttt{rejected}.

The implementation of the **Stage Lock** mechanism and the **Human-in-the-Loop Review Drawer** completely resolved this challenge. Applications marked with \texttt{is\_stage\_locked = true} discard automated mutations, requiring explicit user verification while continuing to log the incoming email to the audit record.

\subsection{Optimizing Neural Models for Low-Memory Serverless Execution}
Hosting large transformer models (such as BERT-Large or LLMs with $>7$ billion parameters) in cloud environments incurs substantial memory footprints ($>8$GB VRAM) and significant financial costs. 

To overcome this, we selected parameter-efficient models (\texttt{bge-small-en-v1.5} and \texttt{gliner\_small-v2.1}) specifically optimized for CPU execution. Quantization and dual-engine fallback architectures ensure that inference executes within 380MB of RAM in under 30ms, completely avoiding costly GPU instances.

\subsection{Mitigating Cloudflare and CloudFront SSL Redirect Loops}
During the production deployment of \texttt{track.mourya.tech}, an \texttt{ERR\_TOO\_MANY\_REDIRECTS} loop occurred when proxying traffic through Cloudflare with the SSL encryption mode set to ``Flexible''. In this mode, Cloudflare contacted the AWS CloudFront edge over unencrypted HTTP (port 80). Because CloudFront was configured with \texttt{ViewerProtocolPolicy: redirect-to-https}, CloudFront issued an HTTP 301 redirect back to HTTPS, which Cloudflare subsequently re-sent over HTTP, causing an infinite loop.

This was resolved by reconfiguring the Cloudflare DNS record to **DNS Only (Grey Cloud)** or enforcing **Full (Strict) SSL**, allowing the AWS Certificate Manager (ACM) SSL certificate on CloudFront to handle end-to-end TLS termination seamlessly.

